\section{First-page contractions and locality}\label{fp:section}

The inverse Poincar\'e pairing on a nearby fiber need not have a useful
expression in terms of surviving classes on the components of a degeneration.
We instead use an inverse pairing on a complex.  Its individual terms have
small support, although only their sum is closed.  The distinction is essential:
we identify a numerical contraction using the closed sum, and subsequently
expand that sum into local terms.

Throughout this section, $\pi:\cY\to\Delta$ is a projective semistable family
of relative dimension $n$, with smooth total space.  Its central fiber
$Y=\bigcup_{v\in V}Y_v$ has smooth components and no triple intersections.
Its dual graph is bipartite.  We orient each edge from color $0$ to color $1$.
An edge $e$ denotes a connected component $D_e$ of a double locus; its two
inclusions are $i_{e,0}$ and $i_{e,1}$.  Disconnected total spaces are allowed.
All cohomology is rational unless complex coefficients or Hodge decompositions
are specified.  The Tate structure $\Q(-1)$ has Hodge type $(1,1)$.
The parity in a complex is its \emph{total} cohomological degree.

\subsection{The three-column complex}

We use the weight spectral sequence with indexing
\begin{equation}\label{fp:eone}
 E_1^{-1,b}=\bigoplus_e H^{b-2}(D_e)(-1),\qquad
 E_1^{0,b}=\bigoplus_v H^b(Y_v),\qquad
 E_1^{1,b}=\bigoplus_e H^b(D_e).
\end{equation}
All other columns are zero.  The differential $d_1$ increases the first index
by one and preserves the second.  Write $C_Y=\operatorname{Tot}(E_1,d_1)$.
There are identifications
\begin{equation}\label{fp:total-complex}
 C_Y^k=\bigoplus_v H^k(Y_v)
 \oplus\bigoplus_e H^{k-1}(D_e)
 \oplus\bigoplus_e H^{k-1}(D_e)(-1).
\end{equation}
Denote the last two kinds of elements by $a$ and $b$, respectively.  We choose
the signs of the residue identifications so that
\begin{equation}\label{fp:differential}
 d(v,a,b)=(\gamma b,\rho v,0),\qquad
 (\rho v)_e=i_{e,1}^*v_1-i_{e,0}^*v_0,
 \qquad
 \gamma b=\sum_e\bigl(-i_{e,0*}b_e+i_{e,1*}b_e\bigr).
\end{equation}
Here a term of $\gamma$ belongs to the indicated component summand.
These conventions are isomorphic to the usual residue conventions by signs
on the summands.

For completeness, $\rho\gamma=0$ can be seen directly.  Distinct double loci
on a fixed component are disjoint.  At $D_e$, the remaining composition is
multiplication by
$c_1(N_{D_e/Y_0})+c_1(N_{D_e/Y_1})$, which is zero because the two normal
bundles are dual in a semistable smoothing.  Thus \eqref{fp:differential}
is a differential.

The Steenbrink theorem identifies
\begin{equation}\label{fp:abutment}
 H^k(C_Y)=\bigoplus_a E_2^{a,k-a}
 =\bigoplus_a\operatorname{Gr}^{W}_{k-a}H^k(Y_t)_{\mathrm{lim}},
\end{equation}
with its pure Hodge structures and the displayed Tate twists; the weight
spectral sequence degenerates at $E_2$
\cite{Steenbrink1976,Steenbrink1977}.
Here $W$ on the right is the usual weight filtration of the limiting mixed
Hodge structure.  In filtered-complex notation below, this incorporates the
usual shift by cohomological degree.

\begin{lemma}[The first-page pairing]\label{fp:pairing}
The complex $C_Y$ has a perfect graded-symmetric chain pairing of degree $2n$
and Tate type $\Q(-n)$.  On elements of total degrees $k$ and $2n-k$ it is
\begin{align}\label{fp:pairing-formula}
 B_Y\bigl((v,a,b),(w,a',b')\bigr)
 ={}&\sum_v\int_{Y_v}v_v\cup w_v
       +(-1)^k\sum_e\int_{D_e}a_e\cup b'_e\\
    &+(-1)^{k+1}\sum_e\int_{D_e}b_e\cup a'_e.\notag
\end{align}
It induces the weight-graded nearby Poincar\'e pairing under
\eqref{fp:abutment}.  Its categorical inverse-pairing tensor $\Delta_{C_Y}$
is closed and is a sum of tensors supported on a single vertex or a single
edge.
\end{lemma}

\begin{proof}
Perfection follows from Poincar\'e duality on the smooth projective $Y_v$ and
$D_e$.  The edge summands pair with one another, not with themselves.  Ordinary
cup-commutativity gives graded symmetry with respect to total degree.
The projection formula says that $\gamma$ is transpose to $\rho$ before the
shift signs.  If $x$ is in a component summand of degree $k$ and $y$ is in a
$b$ summand of degree $2n-k-1$, the two terms of
\begin{equation}\label{fp:chain-pairing}
 B_Y(dx,y)+(-1)^k B_Y(x,dy)=0
\end{equation}
are $(-1)^{k+1}\int \rho x\cup y$ and
$(-1)^k\int \rho x\cup y$.  They cancel.  The case with the $b$ summand
first follows in the same way, and all other cases vanish.  This proves the
chain identity.

We specify the compatibility with nearby integration below, in
\cref{fp:trace-normalization}; it uses the multiplicative comparison and
trace of Fujisawa.  The residue calculation identifying that trace pairing
with \eqref{fp:pairing-formula} is
\cite[Lemma~6.13 and proof of Theorem~8.11]{Fujisawa2014}, after the sign
choices in \eqref{fp:differential}.  Thus this compatibility does not require
choosing representatives for the groups in \eqref{fp:abutment}.

A perfect chain pairing identifies a finite complex with its shifted dual.
Under this identification its inverse-pairing tensor corresponds to the
identity chain map.  It is therefore closed.  Finally,
\eqref{fp:pairing-formula} is block diagonal by vertices and by edges, so its
inverse has the asserted support.  All inverse tensors here and below use
the categorical Koszul convention.
\end{proof}

\begin{lemma}[First Hodge index]\label{fp:first-index}
On $C_Y\otimes\C$, let $P$ act by the first Hodge index, including the Tate
twist.  Then $Pd=dP$.  Consequently $(f(P)\otimes\id)\Delta_{C_Y}$ is closed
for every function $f$ on the finite spectrum of $P$.
\end{lemma}

\begin{proof}
A restriction preserves Hodge type.  A Gysin map raises both Hodge indices
by one, exactly as does the twist in its source $b$ summand.  Thus both arrows
of \eqref{fp:differential} preserve the first index.  The last assertion
follows by applying a chain map to a closed tensor.  Notice that the total
cohomological degree operator does not have this property: $d$ raises total
degree by one.
\end{proof}

\begin{example}\label{fp:curve-example}
Let two rational curves meet in $m$ nodes, where $m=1$ or $2$.  Choose component
units $u_0,u_1$, top classes $t_0,t_1$, and node generators $a_e,b_e$.
Then
\[
 du_0=-\sum_ea_e,\qquad du_1=\sum_ea_e,\qquad
 db_e=-t_0+t_1.
\]
The first Hodge indices of $u,a,b,t$ are $0,0,1,1$.  The inverse tensor is
\[
 \Delta_{C_Y}=\sum_{i=0}^1(u_i\otimes t_i+t_i\otimes u_i)
          +\sum_{e=1}^m(a_e\otimes b_e-b_e\otimes a_e).
\]
Its differential vanishes by cancellation between vertex and edge terms.
For $m=1$ the cohomology dimensions are $(1,0,1)$; for $m=2$ they are
$(1,2,1)$.  Thus the construction includes the graph contribution to nearby
$H^1$.  Individual vertex and edge terms are usually not closed.
\end{example}

\subsection{Filtered multiplicative models and ordered pullbacks}

We use the ordered polynomial-log model of
\cite[Sections~4.14--4.19 and~5.4, Proposition~5.8]{Fujisawa2008}.
In that model each nonempty subset of component indices occurs once, in its
increasing order.  This choice matters: it is not the unrestricted complex
of all injective ordered tuples.  Let $A_Z$ be the Steenbrink complex for a
projective semistable central fiber $Z$, now allowing higher intersections,
and write $K_Z$ for this ordered model.  The standard hypotheses are
properness and smooth K\"ahler components, which hold for our projective
models.  The comparison used below is a bifiltered map
\begin{equation}\label{fp:comparison}
 \phi_Z:A_Z\longrightarrow K_Z
\end{equation}
inducing the limiting mixed-Hodge isomorphism and an isomorphism on second
weight pages.  We give the ordered interpretation of this comparison and
of the trace explicitly, so no identification of different \v{C}ech
conventions is implicit.

Here is the structure of the model that will be used.  If $I$ is a nonempty
set of component indices, $Z_I$ is their intersection with the induced
absolute log structure and $\omega_{Z_I}$ its absolute log de Rham complex.
The local terms are
\[
 \C[u]\otimes\omega_{Z_I},\qquad
 \nabla=1\otimes d+(2\pi\sqrt{-1})^{-1}
                      \partial_u\otimes(d\log s\wedge).
\]
Their component \v{C}ech total complex is $K_Z$.  A power $u^r$ has weight
$2r$ and first Hodge degree $r$.  In \v{C}ech degree $j$ its filtrations are
\begin{align}\label{fp:k-filtrations}
 W_m K_Z&=\sum_{r\geq0}u^r\otimes W_{m+j-2r}\omega_{Z_I},&
 F^pK_Z&=\sum_{r\geq0}u^r\otimes F^{p-r}\omega_{Z_I}.
\end{align}
The $W$ on forms counts logarithmic factors.  Residues describe its graded
pieces by ordinary forms on closed intersections.  The map $\phi_Z$ is a sum
of residue maps multiplied by divided powers of $u$, with universal signs
and Tate factors.  Its summands only involve incident strata
\cite[Definition~5.24 and Lemma~5.25]{Fujisawa2014}.

Multiplication in $K_Z$ is polynomial multiplication, wedge product, and the
ordered \v{C}ech Alexander--Whitney product, with its total-complex signs.
The formulas in \cite[Sections~6.1--6.8]{Fujisawa2014} apply to increasing
tuples: the two overlapping subtuples of an increasing tuple are increasing.
Thus the product is a chain map preserving both filtrations.  Evaluation at
$u=0$ followed by passage to relative log forms respects this product, so its
cohomological product is ordinary nearby cup product.

\begin{lemma}[Ordered comparison]\label{fp:ordered-comparison}
The residue formula for $\phi_Z$, restricted to increasing tuples, gives
\eqref{fp:comparison} with the properties stated above.
\end{lemma}

\begin{proof}
Restriction of a cochain indexed by all injective ordered tuples to increasing
tuples commutes with the \v{C}ech differential, because every face of an
increasing tuple is increasing.  It commutes with the internal differential
and preserves \v{C}ech degree and the two filtrations.  Consequently restricting
the explicit residue map of \cite[Definition~5.24 and Lemma~5.25]{Fujisawa2014}
gives a bifiltered chain map into the ordered model.  This conclusion only
uses the formula and its chain identity, not a quasi-isomorphism assertion
for an unrestricted tuple complex.

Evaluate at $u=0$ and pass to relative log forms.  The comparison square in
\cite[proof of Theorem~5.29]{Fujisawa2014} remains commutative after this
restriction: its equality is an equality of residue maps on each intersection.
Its other arrows are the Steenbrink comparison with relative log de Rham
cohomology, the ordered relative-log \v{C}ech resolution
\cite[Proposition~4.19]{Fujisawa2008}, and the polynomial-log comparison
\cite[Lemma~5.6 and Corollary~5.7]{Fujisawa2008}.  All three induce
isomorphisms on cohomology.  Therefore so does $\phi_Z$.  The rational
comparison is obtained by the same restricted Koszul-residue formula, and
it has the same compatibility with Tate factors as the complex comparison.
The source is a cohomological mixed Hodge complex, and the ordered target
is a weak cohomological mixed Hodge complex by
\cite[Proposition~5.8]{Fujisawa2008}.  Its cohomology carries the mixed
Hodge structure of \cite[Theorem~5.9]{Fujisawa2008}, and its weight spectral
sequence degenerates at $E_2$ by \cite[Lemma~A.6]{Fujisawa2008}.
Thus the cohomological isomorphism is one of mixed Hodge structures.
Weight degeneration and strictness now give
the asserted isomorphism of second pages, exactly as in
\cite[Corollary~5.30]{Fujisawa2014}.
\end{proof}

Let $\widetilde{\cY^{\,r}}$ be the ordered semistable configuration model for
$r$ recorded points, as in \cref{ev:restriction}.  Locally it resolves
\[
 x_i y_i=s,\qquad 1\leq i\leq r,
\]
by the ordered triangulation of $[0,1]^r$.  Its vertices are assignments of
components to the recorded points, and its simplices are chains of assignments
in the coordinatewise order.  Color $0$ precedes color $1$ in each factor.
Order all vertices by the number of color-$1$ coordinates, breaking ties
arbitrarily.  This order is strictly increasing on every nonconstant step in
such a simplex.  Each coordinate projection is therefore weakly increasing
on each ordered simplex.  The same statement applies when some coordinates
are fixed off the double locus, and on disconnected pieces.  Only assignments
lying on a common simplex matter; there is no comparison of points on
disjoint strata.

\begin{lemma}[Ordered pullbacks]\label{fp:pullback}
Each projection of the ordered configuration model to $\cY$ induces a
bifiltered chain pullback of the $K$ models.  On weight pages its terms are
ordinary pullbacks on the relevant intersections and the maps on logarithmic
generators modulo ordinary forms.  These terms are determined by the incident
strata, their normal data, and the chosen ordering.
\end{lemma}

\begin{proof}
Write $f$ for the component assignment of a projection.  For a source
increasing tuple $J=(j_0,\ldots,j_k)$, the component of the pullback of a
\v{C}ech cochain $\eta$ is
\begin{equation}\label{fp:cech-pullback}
 (f^*\eta)_J=
 \begin{cases}
 f_J^*\eta_{f(j_0),\ldots,f(j_k)},&
                 f(j_0)<\cdots<f(j_k),\\
 0,&\text{some projected vertices repeat}.
 \end{cases}
\end{equation}
In the second case the tuple is degenerate.  Formula
\eqref{fp:cech-pullback} commutes with the \v{C}ech differential: if a projected
tuple has exactly one adjacent repetition, deleting either member produces
the same restriction with opposite signs; all other faces remain degenerate.
More repetitions give zero on both sides, and distinct tuples give the usual
pullback identity.  Weak monotonicity ensures that repetitions are adjacent.
This also makes \eqref{fp:cech-pullback} compatible with the
Alexander--Whitney product: a repeated adjacent pair lies in at least one of
the two overlapping subtuples used by that product.

On each intersection use the pullback of absolute log forms and set
$f^*u=u$.  The morphism is over the same base, so $f^*d\log s=d\log s$;
it commutes with $\nabla$.  Ordinary forms pull back to ordinary forms, while
a local log generator pulls back to a sum of log generators and an ordinary
form.  Thus log weight and form degree do not increase.  Every nonzero term
of \eqref{fp:cech-pullback} retains the same \v{C}ech degree, so
\eqref{fp:k-filtrations} proves preservation of both filtrations.

Finally, take residues.  Ordinary corrections to a logarithmic generator
have zero residue, and the remaining maps only use its orders along the
incident divisors.  A unit change of a local boundary equation does not change
these residue maps.  This proves the asserted dependence and locality.
\end{proof}

\subsection{Trace and specialization of correspondences}

If $Z$ has pure dimension $N$, the ordered model has a first-page functional
\begin{equation}\label{fp:theta}
 \Theta_Z:E_1^{0,2N}(K_Z,W)\longrightarrow\C,
 \qquad \Theta_Zd_1=0.
\end{equation}
Here is its precise relation to the residue trace of
\cite[Section~7]{Fujisawa2014}.  Extend an increasing-tuple cochain to all
injective ordered tuples by the sign of the sorting permutation; denote this
map by $\operatorname{Alt}$.  It is a bifiltered chain map, since sorting
intertwines the alternating face sums and leaves internal differentials and
\v{C}ech degree unchanged.  Define $\Theta_Z$ as Fujisawa's explicit
first-page functional composed with $\operatorname{Alt}$.  The formal identity
$\Theta d_1=0$ in \cite[Proposition~7.9]{Fujisawa2014} proves
\eqref{fp:theta}.  This use of the identity does not assume that the
unrestricted injective-tuple complex computes limiting cohomology.

Equivalently, $\Theta_Z$ sums over increasing incidence flags.  For a flag of
length $j+1$ its coefficient is
$(-1)^{j(j-1)/2}(2\pi\sqrt{-1})^j$, followed by the residue after
$d\log s\wedge$ and integration over the closed intersection.  The
$(j+1)!$ orders in the unrestricted formula cancel its factor $1/(j+1)!$:
orientation of the residue and alternating extension have the same sorting
sign.  The functional uses the zero-$u$ part, as in
\cite[Equation~(7.8.1)]{Fujisawa2014}.  In particular it is local on incidence
flags.  Extend it by zero to all other bidegrees of the total first page; it
is then a chain functional.

The same-stratum pairing calculation also holds in this convention.
\cite[Lemma~6.13]{Fujisawa2014} computes the product followed by residue for
one specified ordered tuple.  On that tuple it vanishes except when the two
input residue strata coincide with its underlying intersection and their
indices are opposite.  Restricting that calculation to the increasing tuple
and using the coefficient just given yields the signed integration pairing.
This verifies directly the pairing formula used in \cref{fp:pairing}; it does
not require the alternating-extension map to preserve products.

\begin{lemma}[Normalization of the trace]\label{fp:trace-normalization}
Normalize the residue and Tate factors so that the trace of an extending top
class is the sum of its component integrals.  Under multiplicative nearby
comparison, the descended trace is ordinary integration on every connected
nearby component.  The first-page pairing of \cref{fp:pairing} consequently
induces nearby Poincar\'e duality with this normalization.
\end{lemma}

\begin{proof}
Let $\ell$ be the first Chern class of a relatively ample line bundle.  On a
connected nearby component of dimension $N$, its $N$th power is nonzero and
spans top cohomology.  Its integral is the sum of its integrals over the
components of the corresponding central fiber, by specialization of a top
intersection number.  The residue trace has exactly this value by its stated
normalization.  Hence the two traces agree.  The argument applies separately
to disconnected total families.  After a finite base change, which does not
change the assertion, connected components of smooth fibers can be treated
separately.

Multiplicative comparison identifies cup products.  Applying the identical
traces to cup products identifies the pairings.  Equivalently, the explicit
residue calculation quoted in \cref{fp:pairing} gives component integration
and opposite-column integration on the same double stratum, with no
cross-stratum terms.  This also identifies the induced inverse pairings.
No comparison between two constructions of polarizations on the full limiting
mixed Hodge structure is needed for this argument.
\end{proof}

\begin{lemma}[The leading specialization of an extending class]
\label{fp:leading-class}
Suppose $\widetilde{\cY^{\,r}}\to\Delta$ is smooth and semistable and carries
an algebraic class $\Gamma$ of codimension $c$ restricting to a prescribed
correspondence on the smooth fiber.  The induced weight-graded specialization
of this correspondence is represented on the first page by the ordinary
restrictions $\Gamma|_{Z_v}$ to the smooth components of its central fiber,
inserted through the specialization morphism into $K_Z$.
The same assertion holds with fixed smooth projective source factors and
algebraic correspondences extending over the total family, by first exposing
their cohomology slots and then applying fixed homogeneous inputs.
\end{lemma}

\begin{proof}
The central restriction of $\Gamma$ defines a morphism from $\Q(-c)$ into
central-fiber cohomology, and its image in nearby cohomology is obtained by
the functorial specialization map.  For the ordinary central-fiber
\v{C}ech complex, the weight-$2c$ part of degree $2c$ is represented by
\[
 \bigl(\Gamma|_{Z_v}\bigr)_v\in\bigoplus_v H^{2c}(Z_v).
\]
The restrictions agree on every double intersection, because they come from
one class on the total space.  Thus this tuple is a cycle in the first
ordinary weight differential.  Higher \v{C}ech-degree terms in total degree
$2c$ have smaller weight.

The map from the ordinary \v{C}ech model to the nearby log model sends this
tuple to the corresponding zero-log, zero-$u$ component restrictions.  It is
a morphism of filtered complexes and gives the usual specialization on
cohomology.  Taking its weight-$2c$ graded part gives the asserted
representative.  Strictness of morphisms of mixed Hodge structures ensures
that passing to this graded part loses no part of the induced map from the
pure source.  This statement concerns that induced graded map; it does not
assert that an arbitrary representative in the entire limiting complex has
no lower-weight terms.

For such an extending source correspondence, apply the same argument to
each pure cohomology group of its fixed source, with the prescribed weight
shift.  Exposing its inputs before restriction and applying them afterwards
is the projection formula.  This also treats preintegrated admissible
homogeneous singles.
\end{proof}

For GW theory, the hypothesis that $\Gamma$ is a class on the smooth
\emph{total} configuration family is substantial.  It is furnished by
\cref{ev:restriction}, which constructs the virtual pushforward and computes
its component restrictions before integration.  An arbitrary decomposition
of a cycle on the unresolved special fiber would not suffice for
\cref{fp:leading-class}.

\subsection{Scalar diagrams on complexes}

We record two algebraic facts that explain why the preceding constructions
compute the desired numbers.

\begin{lemma}[Graded limits of scalar diagrams]\label{fp:scalar-limit}
Consider a closed tensor diagram made from flat Hodge tensors of prescribed
types, flat singles of fixed Hodge type, pairings and inverse pairings, and functions
of the first Hodge index on its legs.  Assume its total type is that of a
scalar.  Its value on the smooth fibers is locally constant and equals the
value obtained on the associated Hodge and weight gradeds of the limiting
mixed Hodge structures, with all Tate shifts retained.
\end{lemma}

\begin{proof}
Work in a local flat trivialization.  All factors except the Hodge-index
projectors are constant.  The derivative of a projector onto a summand of the
Hodge decomposition has only off-diagonal blocks: differentiation of
$\Pi_p^2=\Pi_p$ gives $\Pi_p(d\Pi_p)\Pi_p=0$, and differentiation of
$\Pi_p\Pi_q=0$ gives the other diagonal blocks.  In a derivative of the
closed diagram, exactly one such factor has been replaced by an
off-first-index block.  At every other vertex the prescribed Hodge indices
balance.  Summing these balances over the diagram leaves the nonzero index
change at the exceptional factor, so the resulting scalar contraction is
zero.  This proves local constancy; a function on the finite index spectrum
is a linear combination of projectors.

One may equivalently make the entire contraction on the associated Hodge
graded bundles.  The canonical extensions of the Hodge filtrations in a
semistable degeneration are locally free, and flat morphisms of variations
extend as filtered morphisms.  Their graded tensor contraction therefore
extends with the same scalar value.  Finally, the limiting objects are mixed
Hodge structures.  Tensor products, duals, and morphisms are strict for their
weight filtrations.  Taking weight gradeds of the closed contraction, whose
shifts cancel at the scalar output, leaves its value unchanged.  Hodge and
weight gradeds can be taken successively; the resulting double graded spaces
retain the first index including every Tate shift.
\end{proof}

\begin{lemma}[Contraction on a complex]\label{fp:chain-contraction}
Let $C$ be a finite complex with a perfect chain pairing and let $T$ be a
chain functional on a tensor product of copies of $C$, possibly with fixed
cycles and fixed degree shifts.  Contract any paired slots against its closed
inverse-pairing tensor.  The resulting scalar is the contraction of the
induced functional on cohomology against the induced inverse pairing.
The assertion is unchanged if a pairing leg is acted on by a chain map.
\end{lemma}

\begin{proof}
Over a field, the K\"unneth map identifies the cohomology of a finite tensor
product with the tensor product of the cohomologies.  The inverse-pairing
cycle represents the inverse of the induced perfect pairing, since the
chain-level evaluation and coevaluation identities descend to cohomology.
The tensor product of all inserted cycles is a cycle.  A chain functional
annihilates its boundaries, so its value depends only on that cohomology
class.  Applying chain maps to its legs preserves this argument.  Koszul
symmetry gives the same statement for loops: the resulting contraction is
the categorical supertrace.  No splitting of cycles modulo boundaries enters
the calculation.
\end{proof}

\begin{proposition}[Local first-page algebra]\label{fp:local-algebra}
Fix the arities of a scalar test as in \cref{sec:conventions}, applied to
one coefficient operation of \cref{conv:positive-rule}.
For each of its evaluation correspondences, assume the extending class and
assignment-sensitive component restriction rule of \cref{ev:restriction}.
Assume its degree tests extend coherently and its singles are coherent
admissible singles: extending algebraic classes and their complex linear
combinations, or homogeneous cohomology inputs from fixed smooth projective
sources through algebraic correspondences extending over the relevant total
families and inducing flat Hodge maps on their smooth loci.
Then its numerical value has an expansion with these properties.
\begin{enumerate}[label=\textup{(\roman*)}]
\item Each term of its cohomological calculation uses only a bounded number
of positions of the dual graph and closed incidence neighborhoods of bounded
radius.  The bounds depend on the fixed arities and dimensions, not on the
length of inserted chains.
\item The terms are obtained from component restrictions of the evaluation
classes, pullbacks, cups, residues, integrations, and the vertex and edge
blocks of $\Delta_{C_Y}$.  A Hodge-graded coevaluation uses one vertex or one
edge neighborhood.
\item Identical component and stratum diagrams, normal data, orders, and
restriction classes in these neighborhoods give identical terms, with the
same signs and Tate factors.  Changing the smoothing or its connectedness
away from these neighborhoods imposes no extra condition on a term.
\end{enumerate}
Ordinary singles already integrated into an evaluation class need not be
counted as exposed positions.  If an operator subsequently acts on one of
those singles, that slot must instead be exposed and counted.
\end{proposition}

\begin{proof}
First use \cref{fp:scalar-limit} to pass to the Hodge and weight gradeds of
nearby cohomology.  For every coevaluation use the cycle
$\Delta_{C_Y}$ in the small complex \eqref{fp:total-complex}; first-index
factors are chain maps by \cref{fp:first-index}.  Extending cup classes give
chain maps with their usual bidegree shifts.

For a tensor with $r$ exposed evaluation slots, use its own ordered
configuration model.  Transfer each input by $\phi_Y$ and pull it to that
model by \cref{fp:pullback}.  Multiply the inputs in $K_Z$, multiply by the
specialization from \cref{fp:leading-class}, and apply $\Theta_Z$.  All these
operations induce chain maps on first pages, and $\Theta_Z$ is a chain
functional by \eqref{fp:theta}.  The construction represents the desired
cohomological evaluation by multiplicative comparison and
\cref{fp:trace-normalization}.  Use separate models for separate evaluation
tensors; only their exposed slots are linked.  For a classical cup tensor,
use $K_Y$ itself, cup all its inputs there, and apply $\Theta_Y$.  The resulting
closed diagram is therefore a finite tensor contraction of chain functionals
and cycle coevaluations.  By \cref{fp:chain-contraction} it equals the original
weight-graded contraction.  An isomorphism on second pages in
\eqref{fp:comparison} is sufficient: it identifies the induced functionals
and pairings, while all maps used to perform the calculation go forward.

Now expand these actual chain maps and tensors into their residue and
component summands.  A block of the inverse pairing uses one vertex or edge.
A comparison summand uses an incidence of strata.  A product summand uses
intersecting strata in the ordered \v{C}ech diagram.  A trace summand uses
one incidence flag.  Finally, a closed stratum of the ordered configuration
model projects in each target coordinate to one component or to an adjacent
double stratum.  These descriptions prove that each summand is supported in
a union of bounded incidence neighborhoods of the exposed positions.

The bounds are uniform.  A model with $r$ recorded factors has relative
dimension $nr$ and consequently bounded lengths of nonempty incidence flags.
The input weight indices are bounded by the fixed number of legs; by
\eqref{fp:k-filtrations}, only bounded powers of $u$ can occur in the relevant
weight and total degree.  For preintegrated source inputs, these are the
degrees of the image after the Gysin, descendant, and pushforward shifts.
That image lies in the cohomology of the retained configuration, whose
degree range is bounded by its dimension.  Thus the same bound applies
regardless of the number of ordinary auxiliary inputs.
Each slot contributes finitely many local index
types.  Increasing the number of components may increase the number of
summands, but does not enlarge the support or index ranges of any one
summand.  The total number of exposed slots and tensor vertices was fixed in
advance.  This proves (i) and (ii).

Every formula just used consists of ordinary maps on these strata, the
induced normal data, and universal index signs and Tate factors.  Unit changes
in equations disappear after residue.  Hence it is unchanged when the
specified local diagrams are unchanged, proving (iii).  At no step did we
choose a representative of a surviving nearby class or test whether two
positions lie in the same connected smooth component.  Disconnected spaces
are handled by direct sums of their complexes and componentwise traces.
Thus their cohomological connectedness is already encoded by the differential
and its closed Casimir.

Lastly, preintegration changes the coefficient correspondence rather than
its exposed arity.  Its actual extending class and its component restrictions
are still the ones supplied by \cref{ev:restriction}.  Applying a new
cohomological operation to one of its markings requires exposing that slot;
it then contributes one further position.  This gives the final assertion.
\end{proof}

The proposition is a statement about cohomological operations.  The
factorization of an unused collection of relative maps, while retaining the
joint evaluations of other groups, is the separate virtual statement
\cref{ev:groups}.  Together these two statements permit the idle-switch
comparison in \cref{sew:transport}.
